\subsection{Primitive recognition apparatus}
\label{subsec:primitive-recognition}

The sixth entry of the Dirac--Igusa pro-object is the primitive
recognition criterion.  Its conclusion, after the compact source data
are supplied, is an isomorphism of Gram-graded Lie superalgebras
\[
P_X^{\Pi,\mathrm{red}}\cong\mathfrak g_{\Delta_5}
\]
between the cosection-twisted reduced Hall primitive layer of
$K3\times E$ and the imported Borcherds--Kac--Moody denominator
algebra.  Signed superdimensions alone do not suffice for this
theorem.  The Borcherds product fixes
$\sdim P^{\Pi,\mathrm{red}}_{X,\gamma}=f(nm,l)$ for $\gamma=(n,l,m)$;
this is one integer in each degree.  A source--target identification
fixes two integers in every degree (the parity dimensions), the
relation ideal, the radical equality, the no-extra-relations theorem,
the completed PBW filtration, and the Hopf-radical coideal stability.

The determinant admits the zero-bracket graded super vector space with
the same signed dimensions: the denominator is identical, while the
Chevalley diagonal, the non-orthogonal brackets, and the Hopf pairing
datum have disappeared.  Recognition therefore requires
representatives, relations, pairings, kernels, and PBW data, not
numerics.

\subsubsection*{Borcherds--Cartan target datum}

The target side imports the Borcherds--Kac superdatum
$\mathscr B_{\Delta_5}=(H,I,p,\alpha_i,(\,,\,))$ of
Gritsenko--Nikulin \cite{GN,GNPartI,GNII} after the
\cite{KAC:1,BorcherdsGKM88} presentation conventions.  The Cartan is
$H=\Lambda_{II}^{2,1}\otimes\C$, the three real simple roots
$\delta_1,\delta_2,\delta_3$ form a hyperbolic triangle with Gram
matrix
\[
((\delta_i,\delta_j))_{i,j}=
\begin{pmatrix}2&-2&-2\\-2&2&-2\\-2&-2&2\end{pmatrix},
\]
and the imaginary fibres of $I\to H^*$ carry the
Gritsenko--Nikulin parity split
$\mu_{\overline 0}(a),\mu_{\overline 1}(a)$ for imaginary simple
roots.  The degree map is
\[
\alpha(n,l,m)=2nf_2-lf_3+2mf_{-2}
=n\delta_1+m\delta_2+(n+m-l)\delta_3,
\]
with $\widehat{\mathfrak F}(\mathscr B_{\Delta_5})$ the completed
free Lie superalgebra on the parity-homogeneous generators
$e_i,f_i$, $i\in I$, and $J_{\mathrm{BK}}$ the closed
Borcherds--Kac ideal generated by Cartan, Chevalley, real Serre,
orthogonality, and super sign relations.  The target denominator
algebra is
\[
G(\mathscr B_{\Delta_5})
=\widehat{\mathfrak F}(\mathscr B_{\Delta_5})/
\overline{J_{\mathrm{BK}}+\operatorname{Rad}_{\mathrm{GN}}}
=\mathcal A_{\mathrm{den}}
=\mathfrak g_{\Delta_5}.
\]

\subsubsection*{Source-side data}

\begin{theorem}[Primitive recognition from finite source
representatives; \textit{conditional on (S1)--(S10)}]
\label{thm:primitive-recognition}
Let a compact Hall/factorization construction give a completed
$\widehat\Gamma_X$-graded primitive Hall Lie superalgebra
$\widetilde P_X$ and a completed Hopf pairing.  Its normal-ordered
Gram descent
\[
P_X^{\Pi,\mathrm{desc}}
=H^0\Prim_{\mathrm{prot}}(\overline\Pi_{X,*}^\Theta\mathcal F_X)
=\overline\Pi_{X,*}^\Theta\widetilde P_X,
\quad
P_X^{\Pi,\mathrm{red}}
=P_X^{\Pi,\mathrm{desc}}/\overline{\operatorname{Rad}\langle\,,\,\rangle_X}.
\]
Suppose the following ten source-side data are constructed.

\begin{enumerate}
\item[\textup{(S1)}] \emph{Charge and normal-ordered descent.}  The
chain-level descent $\overline\Pi_{X,*}^\Theta$ categorifies the
homomorphism $\overline\Pi_X(c,T)=\Pi_X(c)-T$, so that the descended
bracket is $\Gamma_{\mathrm{gram}}$-graded.  Each descended homogeneous
parity piece is finite dimensional.

\item[\textup{(S2)}] \emph{Cartan and root-degree matching.}  An even
Cartan identification $\iota_H:P^{\Pi,\mathrm{red}}_{X,0}
\xrightarrow{\sim}\Lambda_{II}^{2,1}\otimes\C$ aligning the descended
Gram grading with the GN root grading via $\alpha(n,l,m)$.

\item[\textup{(S3)}] \emph{Simple primitive representatives and
parities.}  For every simple-root index $i\in I$, choose
$\gamma_i\in\Gamma_{\mathrm{gram}}$ with
$\alpha(\gamma_i)=\alpha_i$.  The source supplies parity-homogeneous
representatives $e_i^X,f_i^X,h_i^X$ in
$P^{\Pi,\mathrm{red}}_{X,\gamma_i}$,
$P^{\Pi,\mathrm{red}}_{X,-\gamma_i}$, and
$P^{\Pi,\mathrm{red}}_{X,0}$, with $|e_i^X|=|f_i^X|=p(i)$,
$h_i^X=\iota_H^{-1}(\alpha_i)$, and the GN multiplicity fibres
$\mu_{\overline 0}(a),\mu_{\overline 1}(a)$ for imaginary simple
roots.  Real representatives for $\delta_1,\delta_2,\delta_3$ are even.

\item[\textup{(S4)}] \emph{Borcherds--Kac relations in the Hall
bracket.}  In the protected Hall/factorization bracket the
representatives satisfy:
\[
[h,h']=0,\quad [h,e_i^X]=(h,\alpha_i)e_i^X,\quad
[h,f_i^X]=-(h,\alpha_i)f_i^X,\quad
[e_i^X,f_j^X]=\delta_{ij}h_i^X,
\]
the real Serre relations
\[
(\operatorname{ad}e_i^X)^{1-2(\alpha_i,\alpha_j)/(\alpha_i,\alpha_i)}
e_j^X=0,
\quad
(\operatorname{ad}f_i^X)^{1-2(\alpha_i,\alpha_j)/(\alpha_i,\alpha_i)}
f_j^X=0\qquad(\text{real }i,\ i\ne j),
\]
the Borcherds orthogonality relations
$(\alpha_i,\alpha_j)=0\Rightarrow[e_i^X,e_j^X]=[f_i^X,f_j^X]=0$, and
the super sign rule.  In the rank-three real subsystem, the Serre
exponent is $3$ ($(\delta_i,\delta_j)=-2$); first brackets
$[e_i^X,e_j^X]$ are not forced to vanish.  For the primitive isotropic
degrees $a_{ij}=\delta_i+\delta_j$ ($i<j$), choose
$\gamma_{ij}\in\Gamma_{\mathrm{gram}}$ with
$\alpha(\gamma_{ij})=a_{ij}$.  The source supplies
the nine Gritsenko--Nikulin isotropic simple representatives
$u^{\pm,X}_{ij,r}\in P^{\Pi,\mathrm{red}}_{X,\pm\gamma_{ij}}$,
$1\le r\le 9$, satisfying orthogonality and Chevalley zero with
adjacent real roots, with real-string relations
$(\operatorname{ad}e_k^X)^5 u^{+,X}_{ij,r}=0$ for the complementary
$\{i,j,k\}=\{1,2,3\}$.  The tenth even vector in
$\mathfrak g_{a_{ij}}$ is the real bracket $[e_i^X,e_j^X]$.

\item[\textup{(S5)}] \emph{Counital Hall bialgebra, Hopf pairing, and
radical quotient.}  A continuous completed Hall product, unit,
coproduct, and counit satisfy associativity, coassociativity, the unit
and counit identities, and the bialgebra compatibility
\[
\Delta_X m_X
=(m_X\otimes m_X)(1\otimes\tau\otimes1)(\Delta_X\otimes\Delta_X)
\]
with the Koszul sign braiding \(\tau\).  The bracket is the
supercommutator of this product restricted to primitives.  The Hopf
pairing is continuous, homogeneous, invariant for bracket and coproduct,
has vanishing Frobenius cyclic defect, is non-degenerate after the
radical quotient, and has closed homogeneous radical
$\operatorname{Rad}_X=\{x\mid\langle x,y\rangle_X=0\ \forall y\}$ a Lie
ideal and coideal.  After quotienting by
$\overline{\operatorname{Rad}_X}$, the source radical quotient is
identified with the GN/Kac quotient
$\overline{\operatorname{Rad}_{\mathrm{GN}}}$.  At every finite stage
the Lie-ideal and coideal stability hold and the inverse-limit kernel
is exactly $\operatorname{Rad}_X$.

\item[\textup{(S6)}] \emph{No extra relations.}  If
$\pi:\widehat{\mathfrak F}(\mathscr B_{\Delta_5})\to
P_X^{\Pi,\mathrm{red}}$ is the continuous map determined by the
representatives,
$\ker\pi=\overline{J_{\mathrm{BK}}+\operatorname{Rad}_{\mathrm{GN}}}$.
Equivalently, $G(\mathscr B_{\Delta_5})\to P_X^{\Pi,\mathrm{red}}$ is
injective.  This is checked finite stage by finite stage on saturated
windows, with kernel identifications commuting with HN transitions.

\item[\textup{(S7)}] \emph{Generation.}  $P_X^{\Pi,\mathrm{red}}$ is
topologically generated by $P^{\Pi,\mathrm{red}}_{X,0}$ and the simple
primitive representatives.

\item[\textup{(S8)}] \emph{Completed PBW.}  PBW filtration
compatibility:
\[
\operatorname{gr}_{\mathrm{PBW}}U(P_X^{\Pi,\mathrm{red}})_W
\xrightarrow{\sim}
\operatorname{gr}_{\mathrm{PBW}}U(\mathfrak g_{\Delta_5})_W
\]
on every saturated finite window $W$, with strict transition maps
and inverse limit in the HN topology of
\S\ref{subsec:hall-correspondences}.

\item[\textup{(S9)}] \emph{Exact completion.}  Inverse limit exact on
the saturated finite root windows: transition maps strict, images
closed, and passage to the completed radical quotient commuting with
finite-window kernels, cokernels, PBW associated gradeds, and parity
pieces.  Equivalently, no new relation, generator, radical element,
or parity-cancelling summand appears only after completion.

\item[\textup{(S10)}] \emph{Full parity dimensions.}  For every
nonzero $\gamma$ with $\alpha(\gamma)\in\mathcal R$,
\[
\dim P^{\Pi,\mathrm{red}}_{X,\gamma,\overline p}
=\rootmult_{\overline p}(\alpha(\gamma)),\qquad p=0,1.
\]
The same equality is required in degree \(-\gamma\), with the target
parity dimensions transported by the Borcherds--Kac Chevalley
anti-involution; it is not inferred from the positive half or from the
signed denominator.
On the first timelike channel $\gamma=(1,1,1)$,
\[
d^{\mathrm{GN}}_{(1,1,1),\overline 0}=29,\qquad
d^{\mathrm{GN}}_{(1,1,1),\overline 1}=93.
\]
These are presentation-level target dimensions, not an inference
from $29-93=-64$.
\end{enumerate}
Then $P_X^{\Pi,\mathrm{red}}\cong\mathfrak g_{\Delta_5}$ as completed
Lie superalgebras graded by \(\Gamma_{\mathrm{gram}}\) through
\(\alpha\), with denominator
$64^{-1}\Delta_5(2Z)$.  If in addition the chiral Koszul source datum
and the finite Koszul comparison datum of
\S\ref{subsec:koszul-source} are constructed, including the strict
Koszul Mittag--Leffler exactness for the compatibility complexes, then
\(\Omega_E^{\mathrm{ch}}C_X\simeq\mathsf{Den}_{\Delta_5,E}\) through a
map respecting the Weyl action, Pfaffian orientation character, Hall
pairing, radical quotient, and PBW associated graded.
\end{theorem}

\begin{proof}
The simple primitive representatives define a continuous homomorphism
$\pi:\widehat{\mathfrak F}(\mathscr B_{\Delta_5})\to
P_X^{\Pi,\mathrm{red}}$.  (S1) and (S2) make $\pi$ a homomorphism of
the descended Gram-graded completions.  (S3) fixes the real and
imaginary simple generators with their parities and multiplicity
fibres.  (S4) puts $J_{\mathrm{BK}}$ in the kernel.  (S5) identifies
the compact Hopf radical with the GN/Kac radical, so $\pi$ descends
to $G(\mathscr B_{\Delta_5})\to P_X^{\Pi,\mathrm{red}}$.  (S7) gives
surjectivity; (S6) gives injectivity after the completed radical
quotient.  (S8) gives PBW-compatible finite-height identifications;
(S9), in the form of
Proposition~\ref{prop:strict-ml-completion-primitive-recognition},
gives exactness in the inverse limit; (S10) rules out hidden
zero-superdimension root spaces and cancelling ideals invisible to the
determinant.  The chiral statement uses $C_X$, $b_X$,
$\Theta_{\mathrm{Kos}}$, the finite Koszul comparison datum
\(\mathfrak K^{\mathrm{cmp}}_{X,R}\), and the Koszul
Mittag--Leffler exactness of \S\ref{subsec:koszul-source}.
\end{proof}

\subsubsection*{Cofinal finite-window primitive recognition}

The Borcherds product fixes signed target superdimensions; the
Gritsenko--Nikulin/Kac presentation fixes full target parity
dimensions only in the degrees where the presentation has been
computed.  A compact $K3\times E$ source is identified with the
denominator algebra only by source representatives, pairings,
radicals, relation ideals, PBW, and exact transition maps.

The imported target on the first window $W_{\le 3}$ has parity table
\[
\delta_i:1|0,\quad
a_{ij}:10|0,\quad
2\delta_i+\delta_j:1|0,\quad
\delta_{123}:29|93,
\]
so $29-93=-64=f(1,1)$.  The row $a_{ij}:10|0$ is the real bracket
$[e_i,e_j]$ together with the nine Gritsenko--Nikulin isotropic simple
directions.  These are target data; they do not construct
compact source primitives, source pairings, radical kernels, or
source bracket constants.
The finite target packet \(\texttt{wle3\_target\_parity}\) verifies
this table from the theta-series coefficients and the
Gritsenko--Nikulin/Kac presentation: \(a_{ij}\) is isotropic,
the real-string exponent for \(2\delta_i+\delta_j\) is \(3\), and
\(\delta_{123}\) splits as \(29=2+27\) even directions and \(93\) odd
imaginary-simple directions.  It is a target-reference table only.

On the first base terminal and saturation-extension target rows the
coefficient table
\(\texttt{k3e\_first\_relation\_window\_coefficients}\) computes
\[
f(4,4)=10,\quad f(2,2)=108,\quad
f(2,1)=f(4,3)=-513,\quad
f(3,3)=-64,\quad f(4,2)=4016,
\]
together with their discriminant-orbit translates and the terminal
zero rows \(f(5,5)=f(1,-3)=0\).  These equalities match the signed
dimensions in the A071 target presentation.  They remain one-integer
target arithmetic; they do not supply the two parity dimensions required
by (R3), or any compact-source row in (R0)--(R8).
The A071 target-presentation check verifies the target rows
\[
2a_{ij}:10|0,\qquad C_{k,3}:29|93,\qquad
C_{k,4}:10|0,\qquad C_{k,5}:0|0,
\]
and verifies that \(C_{k,2}\), \(D_k=\delta_k+2a_{ij}\), and
\(2\delta_{123}\) are signed-only: they have no target basis,
positive-negative pairing block, radical block, or PBW row.  This is a
target-reference computation, not a compact-source theorem.

The first relation-closed compact-source window is therefore a separate
finite computation.  Its row packet is
\[
\mathrm{window\_closure},\quad
\mathrm{target\_source\_representatives},\quad
\mathrm{chevalley\_serre\_matrices},\quad
\mathrm{hall\_boundary\_complex},\quad
\mathrm{spectral\_sequence},\quad
\mathrm{kernel\_matrices},\quad
\mathrm{pbw\_graded},\quad
\mathrm{first\_window\_theorems},\quad
\mathrm{transitions},\quad
\mathrm{scalar\_firewall}.
\]
These rows must prove downward saturation and relation closure,
geometric source representatives with full parity ranks, Cartan,
Chevalley, real-Serre, Borcherds-orthogonality, and super-sign
relations, the Hall--Chevalley boundary \(d_1\), strong convergence of
the finite-window spectral sequence, computed kernel bases, no-extra
kernel equality, PBW associated-graded equality, and strict
Mittag--Leffler transition.  Target labels, signed dimensions, scalar
traces, Pfaffian products, denominator products, target PBW rows,
arbitrary matrices, and status-only rows are excluded by the scalar
firewall.

\begin{proposition}[Target-presentation rows are a prerequisite;
\textit{proved}]
\label{prop:target-presentation-prerequisite}
Let \(W\subset\mathcal R_+\) be a finite downward-saturated target-root
window.  If \(W\) is used in a primitive-recognition datum satisfying
(R0)--(R8), then the target restriction
\(\mathfrak g_{\Delta_5}|_{S_W}\) must be represented by a finite
target-presentation datum in the sense of
Definition~\ref{def:bkm-finite-target-presentation-datum}.  Signed
rows \(\smult(\beta)\) and imaginary-simple integers \(m(\beta)\)
alone are not admissible codomains for the comparison maps
\(A_{\beta,\overline p}\).
\end{proposition}

\begin{proof}
Clause (R2) compares source relation rows with the Borcherds--Kac
target relations, so the Cartan, Chevalley, real-Serre,
orthogonality, and super-sign rows must be present in the target
window.  Clause (R3) compares two parity dimensions in each sign, not
the signed difference.  Clause (R4) compares the source
positive-negative Hopf pairing and radical quotient with the target
invariant-pairing kernel.  Clause (R7) compares PBW associated gradeds.
Thus the target side must carry the parity, relation, pairing,
radical, and PBW rows of
Definition~\ref{def:bkm-finite-target-presentation-datum}.  A scalar
row \(\smult(\beta)\) or \(m(\beta)\) supplies none of the pairing,
radical, relation, or PBW codomains, and supplies only one integer
where (R3) requires two.
\end{proof}

\begin{definition}[Cofinal finite-window primitive-recognition datum]
\label{def:cofinal-primitive-recognition-datum}
A finite set \(W\subset\mathcal R_+\) is \emph{downward saturated} if
\(\beta\in W\), \(\beta'\in\mathcal R_+\), \(0<\beta'\le\beta\), and
\(\beta-\beta'\in Q_+\) imply \(\beta'\in W\).  The canonical
height-window convention is
\[
W_\nu=\{\beta=b_1\delta_1+b_2\delta_2+b_3\delta_3\in\mathcal R_+
\mid b_1+b_2+b_3\le \nu\}.
\]
A cofinal finite-window primitive-recognition datum for
$P_X^{\Pi,\mathrm{red}}$ is an increasing sequence of finite
downward-saturated target-root windows
$W_1\subset W_2\subset\cdots$ with
$\bigcup_\nu W_\nu=\mathcal R_+$, together with the following finite
data on every $W_\nu$, compatible under $W_\nu\subset W_{\nu+1}$.  Put
\[
\Gamma(W_\nu):=\{\gamma\in\Gamma_{\mathrm{gram}}\mid
\alpha(\gamma)\in W_\nu\}.
\]

\begin{enumerate}
\item[(R0)] \emph{Compact source labels.}  Every source basis vector
of degree in $\Gamma(W_\nu)\cup(-\Gamma(W_\nu))\cup\{0\}$ satisfies
the compact source admissibility condition.  Target labels from
$G(\mathscr B_{\Delta_5})$ appear only in target rows or as
codomains of comparison maps.
\item[(R1)] \emph{Representatives and parities.}  Simple-root
representatives $e_i^X,f_i^X,h_i^X$ whose target roots lie in $W_\nu$
are supplied in source degrees $\pm\gamma_i$ with
$\alpha(\gamma_i)=\alpha_i$, with the GN parity fibres for imaginary
simple roots and even real representatives for $\delta_1,\delta_2,\delta_3$.
\item[(R2)] \emph{Relations.}  Cartan, Chevalley, real Serre,
orthogonality, and super sign relations checked for every relation
whose displayed target roots lie in
$W_\nu\cup(-W_\nu)\cup\{0\}$ and whose source lifts lie in
$\Gamma(W_\nu)\cup(-\Gamma(W_\nu))\cup\{0\}$.
\item[(R3)] \emph{Full finite parity dimensions.}  For every
$\gamma\in\Gamma(W_\nu)$ and for its negative \(-\gamma\),
$\dim P^{\Pi,\mathrm{red}}_{X,\gamma,\overline p}
=\rootmult_{\overline p}^{\mathrm{GN}}(\alpha(\gamma))$ ($p=0,1$),
with the negative-root parity dimensions carried by the
Borcherds--Kac Chevalley anti-involution, and with zero on source
degrees whose image is not a target root.
Two-integer equality, not signed-superdimension.
\item[(R4)] \emph{Hall bialgebra, Hopf pairing, and radical.}  The
finite product \(M\), coproduct \(D\), unit \(\eta\), and counit
\(\epsilon\) satisfy associativity, coassociativity, unit, counit, and
bialgebra compatibility identities on the window; the supercommutator
bracket sends primitives to primitives.  Positive-negative pairing
matrices are written in parity blocks; kernels agree with GN/Kac
invariant-pairing kernels; Hopf adjointness, the primitive Frobenius
cyclic identity, and non-degeneracy on the radical quotient are verified
by finite rank-zero defect rows; stable under finite brackets in the
window; coideal for finite coproduct; carried to the next window by
transitions; inverse-limit radical the closure of finite kernels.
\item[(R5)] \emph{No extra relations.}  Finite kernel equality
$\ker\pi_\nu=
(\overline{J_{\mathrm{BK}}+\operatorname{Rad}_{\mathrm{GN}}})_{\le W_\nu}$.
\item[(R6)] \emph{Generation.}  Source primitive spaces over
$\Gamma(W_\nu)\cup(-\Gamma(W_\nu))$ are generated by Cartan and
supplied simple representatives through brackets whose intermediate
target roots stay in $W_\nu$.
\item[(R7)] \emph{Completed PBW compatibility.}  Strict PBW
filtration preservation under transitions, inverse-limit PBW
compatibility.
\item[(R8)] \emph{Exact triangular completion.}  For any source stage
\(R\) assigned to \(W_\nu\), words whose total Gram degree lies in the
finite target test window \(\Gamma_R^{\mathrm{test}}\) of
Definition~\ref{def:finite-target-test-window} only through
positive-negative cancellation are not added as extra finite generators;
transition maps strict, images closed, \(\lim^1\) vanishing on the
triangular finite-window systems.
\end{enumerate}
The datum is finite at each $\nu$, but not numerical: it contains
representatives, relation verifications, pairing matrices, kernel
equalities, and PBW transition data.  The Borcherds product and the
OP scalar square provide none of these source-side objects.  The further
assertion that every relation-closed target degree occurs at some finite
HN stage is the finite-\(R\) exhaustion datum of
Definition~\ref{def:relation-closed-target-finite-r-exhaustion}; it is
not a consequence of target relation closure alone.
\end{definition}

\begin{proposition}[Height root windows and their Gram preimages;
\textit{proved}]
\label{prop:height-root-windows-gamma}
The sets \(W_\nu\) of
Definition~\ref{def:cofinal-primitive-recognition-datum} are finite,
downward saturated, and cofinal in \(\mathcal R_+\).  If
\(\beta=b_1\delta_1+b_2\delta_2+b_3\delta_3\in W_\nu\), then
\[
\gamma_\beta=(b_1,\ b_1+b_2-b_3,\ b_2)\in\Gamma_{\mathrm{gram}}
\]
is the unique element with
\[
\alpha(\gamma_\beta)
=b_1\delta_1+b_2\delta_2+b_3\delta_3.
\]
Hence
\[
\Gamma(W_\nu)
=\{(b_1,b_1+b_2-b_3,b_2)\mid
b_1\delta_1+b_2\delta_2+b_3\delta_3\in W_\nu\}.
\]
Moreover \(\gamma_\beta\) lies in the type-II product chamber:
\[
b_2>0\Rightarrow m>0,\qquad
b_2=0,\ b_1>0\Rightarrow m=0,\ n>0,\qquad
b_1=b_2=0,\ b_3>0\Rightarrow m=n=0,\ l<0.
\]
\end{proposition}

\begin{proof}
Finiteness follows because \(b_1+b_2+b_3\le\nu\) gives only finitely
many triples in \(Q_+\).  If
\(\beta'\in\mathcal R_+\), \(0<\beta'\le\beta\), and
\(\beta\in W_\nu\), then
\(\operatorname{ht}(\beta')\le\operatorname{ht}(\beta)\le\nu\); hence
\(\beta'\in W_\nu\).  Cofinality follows because every positive root
has finite height.

For \(\gamma=(n,l,m)\), the product-chamber root map is
\[
\alpha(n,l,m)=n\delta_1+m\delta_2+(n+m-l)\delta_3.
\]
Solving \(\alpha(n,l,m)=b_1\delta_1+b_2\delta_2+b_3\delta_3\) gives
\(n=b_1\), \(m=b_2\), and \(l=b_1+b_2-b_3\), proving uniqueness and the
displayed formula for \(\Gamma(W_\nu)\).  The three chamber cases are
then immediate from \(b_i\ge0\) and \(\beta\ne0\).
\end{proof}

The packet
\(\texttt{certificates/lattice/downward\_root\_windows}\) records
Proposition~\ref{prop:height-root-windows-gamma}.  It verifies
\texttt{DOWNWARD\_ROOT\_WINDOWS\_VERIFIED}: the height-window formula,
the inverse formula for \(\Gamma(W_\nu)\), the product-chamber sector
rule, and the ambient positive-cone audit for \(1\le\nu\le7\).  It is a
target-side root-window convention, not a compact-source recognition
packet, not a Pfaffian orientation, and not a protected trace.

\begin{proposition}[Active height windows exhaust the positive roots;
\textit{proved}]
\label{prop:active-root-window-exhaustion}
Let \(W_\nu^{\mathrm{act}}:=W_\nu\).  Then
\[
\bigcup_{\nu\ge1}W_\nu^{\mathrm{act}}=\mathcal R_+ .
\]
Equivalently, every positive root
\(\beta=b_1\delta_1+b_2\delta_2+b_3\delta_3\) is witnessed in the
finite active window \(W_{\operatorname{ht}(\beta)}\), where
\(\operatorname{ht}(\beta)=b_1+b_2+b_3\).
\end{proposition}

\begin{proof}
By definition \(W_\nu\subset\mathcal R_+\) for every \(\nu\).  Conversely,
if \(\beta\in\mathcal R_+\), then its coefficients in the positive
root monoid are non-negative and have finite sum
\(\operatorname{ht}(\beta)\).  Hence
\(\beta\in W_{\operatorname{ht}(\beta)}^{\mathrm{act}}\).  The two
inclusions give the displayed equality.
\end{proof}

The packet
\(\texttt{certificates/lattice/active\_root\_window\_exhaustion}\)
records Proposition~\ref{prop:active-root-window-exhaustion}.  It
verifies \texttt{ACTIVE\_ROOT\_WINDOW\_EXHAUSTION\_VERIFIED}: the
height-witness formula
\(\beta\in\mathcal R_+\Rightarrow\beta\in W_{\operatorname{ht}(\beta)}\),
the low active membership samples, the inverse \(\gamma_\beta\) formula,
and the type-II chamber sector rule.  It is not the compact-source
support exhaustion required before moduli transitions, orientation
transport, vanishing-cycle transport, Hall operations, radical
transport, or PBW transport.

\begin{definition}[Source-window compact-support exhaustion datum]
\label{def:source-window-compact-support-exhaustion}
A source-window compact-support exhaustion datum consists of:
\begin{enumerate}
\item a set \(\mathscr S_X^{\mathrm{cpt}}\) of compact-source support
rows, each carrying a source degree, a compact geometric provenance row,
proper support, finite stabilizer data, and the orientation row used by
the source Hall construction;
\item a height function
\[
h_X:\mathscr S_X^{\mathrm{cpt}}\longrightarrow \mathbb Z_{\ge0}
\]
with finite fibres;
\item finite source windows
\[
\mathscr S_{X,\le R}^{\mathrm{cpt}}
=\{s\in\mathscr S_X^{\mathrm{cpt}}\mid h_X(s)\le R\};
\]
\item a cofinality witness proving that every
\(s\in\mathscr S_X^{\mathrm{cpt}}\) lies in
\(\mathscr S_{X,\le h_X(s)}^{\mathrm{cpt}}\);
\item transition rows
\(\mathscr S_{X,\le R}^{\mathrm{cpt}}\hookrightarrow
\mathscr S_{X,\le R'}^{\mathrm{cpt}}\) for \(R\le R'\), later refined
to closed embeddings or proper maps before orientation, vanishing-cycle,
Hall, radical, and PBW transport is asserted.
\end{enumerate}
Target root labels, signed multiplicities, Pfaffian products, and
protected traces are not entries of this datum.
\end{definition}

\begin{proposition}[Source-window compact-support exhaustion criterion;
\textit{proved}]
\label{prop:source-window-compact-support-criterion}
If a datum of
Definition~\ref{def:source-window-compact-support-exhaustion} is
supplied, then the finite source windows exhaust compact support:
\[
\bigcup_{R\ge0}\mathscr S_{X,\le R}^{\mathrm{cpt}}
=\mathscr S_X^{\mathrm{cpt}}.
\]
\end{proposition}

\begin{proof}
The inclusion from left to right is part of the definition of the
finite source windows.  Conversely, for
\(s\in\mathscr S_X^{\mathrm{cpt}}\), the cofinality witness gives
\(s\in\mathscr S_{X,\le h_X(s)}^{\mathrm{cpt}}\).  Thus every compact
source support row occurs in one finite source window.
\end{proof}

The packet
\(\texttt{certificates/sources/source\_window\_compact\_support}\)
records the present state of this datum.  Its verifier status is
\texttt{SOURCE\_WINDOW\_COMPACT\_SUPPORT\_OBSTRUCTION\_VERIFIED}: the
obstruction ledger lists the missing compact support index set, source
height function, finite source windows, compact provenance rows,
cofinality witness, source-window transitions, and target-source
compatibility rows.  Its core source tables are empty, so
the compact-source support-exhaustion theorem remains open.  This is
the firewall separating the target-root exhaustion of
Proposition~\ref{prop:active-root-window-exhaustion} from compact-source
support exhaustion.

\begin{definition}[Target/source window compatibility datum]
\label{def:target-source-window-compatibility}
Let \(W_\nu\) be the target height window and
\(\mathscr S_{X,\le R}^{\mathrm{cpt}}\) a finite compact-source
window.  A target/source window compatibility datum consists of:
\begin{enumerate}
\item a degree map
\[
d_R:\mathscr S_{X,\le R}^{\mathrm{cpt}}\longrightarrow
\Gamma_{\mathrm{gram}}
\]
recording the normal-ordered Gram degree of every compact support row;
\item a target index \(\nu(R)\) such that
\[
d_R(s)\in\Gamma(W_{\nu(R)})\qquad
\text{for every }s\in\mathscr S_{X,\le R}^{\mathrm{cpt}};
\]
\item the row-wise equality
\[
\alpha(d_R(s))\in W_{\nu(R)}
\]
with zero degree-map defect;
\item compatibility with compact provenance: the support row, its
source degree, its compact geometric provenance, and its Gram degree
all carry the same support identifier;
\item compatibility under transitions
\[
\mathscr S_{X,\le R}^{\mathrm{cpt}}\hookrightarrow
\mathscr S_{X,\le R'}^{\mathrm{cpt}},\qquad
W_{\nu(R)}\subset W_{\nu(R')},
\]
so that the source degree map commutes with passage to larger windows.
\end{enumerate}
\end{definition}

\begin{proposition}[Compatibility criterion for target and source windows;
\textit{proved}]
\label{prop:target-source-window-compatibility-criterion}
Given a datum of
Definition~\ref{def:target-source-window-compatibility}, every compact
source support row in \(\mathscr S_{X,\le R}^{\mathrm{cpt}}\) has a
well-defined target codomain in
\(\Gamma(W_{\nu(R)})\cup(-\Gamma(W_{\nu(R)}))\cup\{0\}\), and the
finite comparison maps on that source window are degree-compatible with
the target window \(W_{\nu(R)}\).
\end{proposition}

\begin{proof}
For a positive source degree the codomain is \(\alpha(d_R(s))\in
W_{\nu(R)}\) by the compatibility datum.  The negative degree is carried
to \(-\Gamma(W_{\nu(R)})\), and the Cartan degree is \(0\).  The
zero-defect equality makes the target degree independent of notation,
and the transition condition shows that enlarging \(R\) does not move a
previously defined source support row outside the enlarged target
window.  Thus every finite comparison row has a target codomain in the
chosen window.
\end{proof}

The packet
\(\texttt{certificates/sources/target\_source\_window\_compatibility}\)
records the current state of this datum.  Its verifier status is
\texttt{TARGET\_SOURCE\_WINDOW\_COMPATIBILITY\_OBSTRUCTION\_VERIFIED}:
the source degree maps, \(\Gamma(W_\nu)\)-membership rows, target-window
matching rows, compact-provenance compatibility rows, cofinal-index
compatibility rows, transition compatibility rows, and zero
compatibility-defect rows are not yet supplied.  Hence target/source
window compatibility remains open.

\begin{proposition}[Finite-window presentation recognition;
\textit{proved}]
\label{prop:finite-window-presentation-recognition}
Fix a downward-saturated finite target-root window \(W\) and put
\(S_W=W\cup(-W)\cup\{0\}\).  For a graded Lie superalgebra \(L\), write
\[
L|_{S_W}:=\bigoplus_{\beta\in S_W}L_\beta
\]
with the bracket retained only on pairs whose total degree remains in
\(S_W\).  Suppose a compact \(K3\times E\) source supplies
\((\mathrm{R0})\)--\((\mathrm{R8})\) on \(W\).  Then the finite
comparison matrices \(A_{\beta,\bar p}\) determine an isomorphism of
truncated Gram-graded Lie superalgebras
\[
A_W:\ P^{\Pi,\mathrm{red}}_X|_{\alpha^{-1}(S_W)}
\xrightarrow{\ \sim\ }
\mathfrak g_{\Delta_5}|_{S_W},
\]
compatible with the Cartan identification, the simple-root
representatives, the Hopf-pairing radical quotient, and the
associated-graded PBW filtration on \(W\).  Conversely, such an
isomorphism, together with compact-source basis provenance, parity
blocks, simple representatives, product, coproduct, unit, counit,
Hopf-pairing, radical, quotient, and transition maps from which its
matrices are obtained, gives \((\mathrm{R0})\)--\((\mathrm{R8})\).
\end{proposition}

\begin{proof}
The representatives of (R1) define a homomorphism from the free
Borcherds--Kac presentation to the source truncation.  The relation
rows of (R2) put every Cartan, Chevalley, real-Serre, orthogonality,
and super-sign relation whose degree lies in \(S_W\) in the kernel.
The Hopf bialgebra and pairing data of (R4) make the source quotient a
truncated Lie superalgebra: product and coproduct are compatible, the
supercommutator closes on primitives, and the radical is a Lie ideal
and a coproduct coideal.  The no-extra relation row (R5) is exactly the
finite equality
\[
\ker\pi_W=
\bigl(\overline{J_{\mathrm{BK}}+\operatorname{Rad}_{\mathrm{GN}}}
\bigr)|_{S_W}.
\]
Thus the presentation map descends injectively from
\(\mathfrak g_{\Delta_5}|_{S_W}\) to the source radical quotient.
Generation (R6) gives surjectivity on every source degree in
\(\alpha^{-1}(S_W)\).  The parity dimensions of (R3) identify the
even and odd dimensions degree by degree in both signs, so no
zero-superdimension summand or parity-cancelling kernel can remain.
PBW compatibility (R7) identifies the associated gradeds of the
universal enveloping algebras on \(W\).  Exact triangular completion
(R8) says that the finite truncation is stable under passage to the
next window and that no relation is created only by completion.  The
matrix \(A_W\) is therefore a degree-preserving bijection preserving
all brackets, pairing quotients, and PBW associated gradeds present in
the finite window.

Conversely, a truncated isomorphism with specified compact-source basis
provenance, parity blocks, simple representatives, product, coproduct,
unit, counit, Hopf pairing, radical, quotient, and transition maps
yields the relation rows, radical identities, no-extra equality,
generation spans, PBW rows, and transition rows by writing those finite
maps in homogeneous bases.  Without the named compact-source maps one
has only a linear isomorphism of graded vector spaces, not (R0)--(R8).
\end{proof}

\begin{definition}[PBW-transition datum]
\label{def:transition-pbw-filtration-datum}
Let \(W_\nu\subset W_{\nu+1}\) be two downward-saturated finite
windows.  Put
\[
L_\nu=P^{\Pi,\mathrm{red}}_{X,\nu}/\operatorname{Rad}^{\Pi}_{X,\nu},
\qquad
L_{\nu+1}=P^{\Pi,\mathrm{red}}_{X,\nu+1}/
\operatorname{Rad}^{\Pi}_{X,\nu+1}.
\]
A PBW-transition datum for \(L_{\nu+1}\to L_\nu\) consists of the
following finite rows:
\begin{enumerate}
\item radical-quotient transition matrices \(T^Q_{\nu+1,\nu}\) in
homogeneous quotient bases;
\item ordered PBW word bases for \(F^{\mathrm{PBW}}_{\le d}U(L_\nu)\)
and \(F^{\mathrm{PBW}}_{\le d}U(L_{\nu+1})\), with their
associated-graded ranks;
\item source and target relation-rewrite rows which identify the
normal forms in the enveloping algebras;
\item the multiplicative word transition induced by
\(T^Q_{\nu+1,\nu}\), together with the zero-defect relation-rewrite
identity saying that it descends to
\[
U(T^Q_{\nu+1,\nu}):U(L_{\nu+1})\longrightarrow U(L_\nu);
\]
\item for every retained \(d\), a zero-defect filtered-image row
\[
U(T^Q_{\nu+1,\nu})
\bigl(F^{\mathrm{PBW}}_{\le d}U(L_{\nu+1})\bigr)
\subset
F^{\mathrm{PBW}}_{\le d}U(L_\nu);
\]
\item the induced matrices on
\(\operatorname{gr}^{\mathrm{PBW}}_dU(L_{\nu+1})\to
\operatorname{gr}^{\mathrm{PBW}}_dU(L_\nu)\), the PBW-symbol
intertwining identity, the \(A_\beta\)-PBW transition square, and the
composition law for three consecutive windows.
\end{enumerate}
The datum is absent if only primitive transitions, radical transitions,
PBW rank equalities, target PBW rows, Hilbert series, scalar traces,
or the formal normal-ordered PBW-pushforward packet are supplied.
\end{definition}

\begin{proposition}[Transition preservation of PBW filtrations;
\textit{proved}]
\label{prop:transition-pbw-filtration-preservation-criterion}
A PBW-transition datum in the sense of
Definition~\ref{def:transition-pbw-filtration-datum} determines a
filtered algebra homomorphism
\[
U(T^Q_{\nu+1,\nu}):
\bigl(U(L_{\nu+1}),F^{\mathrm{PBW}}_{\le\bullet}\bigr)
\longrightarrow
\bigl(U(L_\nu),F^{\mathrm{PBW}}_{\le\bullet}\bigr)
\]
and hence maps
\[
\operatorname{gr}^{\mathrm{PBW}}_d U(L_{\nu+1})
\longrightarrow
\operatorname{gr}^{\mathrm{PBW}}_d U(L_\nu)
\]
for every retained PBW degree \(d\).  These associated-graded maps are
compatible with the finite comparison matrices \(A_\beta\), and they
compose functorially on triples of windows.  Without the rows of
Definition~\ref{def:transition-pbw-filtration-datum}, transition
preservation of PBW filtrations is not established.
\end{proposition}

\begin{proof}
The quotient transition matrix sends degree-one primitive generators of
\(L_{\nu+1}\) to degree-one elements of \(L_\nu\).  Multiplication gives
a word-level map from the tensor algebra on \(L_{\nu+1}\) to the tensor
algebra on \(L_\nu\).  The relation-rewrite zero-defect row is exactly
the statement that the image of every source enveloping relation is a
target enveloping relation; hence the word-level map descends to the
universal enveloping quotients.

The filtered-image row says, in the chosen finite PBW bases, that the
matrix of the descended map has no component from words of length
\(\le d\) to target normal forms of length \(>d\).  This is precisely
\(U(T^Q_{\nu+1,\nu})(F^{\mathrm{PBW}}_{\le d})\subset
F^{\mathrm{PBW}}_{\le d}\).  Therefore the standard quotient
\(F_{\le d}/F_{\le d-1}\) produces the displayed associated-graded
matrix.  The PBW-symbol row identifies this quotient map with the
symbol of the degree-one transition, and the \(A_\beta\)-square is the
same matrix identity after applying the finite source-to-target
comparison.  The composition row gives equality of the two matrix
products for three windows, so the associated graded maps form an
inverse system.  No scalar or rank-only datum contains these matrix
identities.
\end{proof}

The packet
\[
\texttt{certificates/hall/transition\_pbw\_filtration\_preservation}
\]
records the current state of this datum.  Its verifier returns
\(\texttt{TRANSITION\_PBW\_FILTRATION\_PRESERVATION\_OBSTRUCTION\_VERIFIED}\):
the radical-quotient transition input, source and target PBW
filtration rows, associated-graded rank rows, ordered word bases,
relation-rewrite systems, quotient transition matrices, multiplicative
word transitions, relation-rewrite compatibility rows, filtered-image
identities, restricted filtered transition matrices, associated-graded
transition matrices, PBW-symbol identities, \(A_\beta\)-PBW transition
squares, and PBW-transition composition rows are not yet supplied.

\begin{definition}[Parity-transition datum]
\label{def:transition-parity-decomposition-datum}
Let \(V_{\nu,\gamma}\) and \(V_{\nu+1,\gamma}\) be the retained
finite primitive source spaces, after the finite radical quotient if
recognition is being tested after \textup{(R4)}.  A parity-transition
datum for
\[
T_{\nu+1,\nu,\gamma}:V_{\nu+1,\gamma}\longrightarrow V_{\nu,\gamma}
\]
consists of parity-homogeneous bases
\[
V_{\nu,\gamma}=V_{\nu,\gamma,\overline 0}\oplus
V_{\nu,\gamma,\overline 1},
\qquad
V_{\nu+1,\gamma}=V_{\nu+1,\gamma,\overline 0}\oplus
V_{\nu+1,\gamma,\overline 1},
\]
the finite fermion-parity involutions
\[
J_{\nu,\gamma}=+1\ \text{on }V_{\nu,\gamma,\overline 0},\qquad
J_{\nu,\gamma}=-1\ \text{on }V_{\nu,\gamma,\overline 1},
\]
and the transition matrix of \(T_{\nu+1,\nu,\gamma}\), with the
following zero-defect rows:
\[
J_{\nu,\gamma}T_{\nu+1,\nu,\gamma}
=T_{\nu+1,\nu,\gamma}J_{\nu+1,\gamma},
\]
equivalently
\[
T_{\overline 0\overline 1}=0,\qquad
T_{\overline 1\overline 0}=0.
\]
The datum also includes the restricted matrices
\[
T_{\overline p}:V_{\nu+1,\gamma,\overline p}
\longrightarrow V_{\nu,\gamma,\overline p},
\qquad p=0,1,
\]
the same rows for \(-\gamma\) transported by the Borcherds--Kac
Chevalley anti-involution, the comparison square with the finite
\(A_\beta\)-matrices on each parity piece, and the composition law for
three consecutive windows.  Signed superdimension, target parity
tables, parity ranks without transition matrices, primitive/radical/PBW
transition rows, and the formal parity-pushforward packet are not such
a datum.
\end{definition}

\begin{proposition}[Transition preservation of parity decompositions;
\textit{proved}]
\label{prop:transition-parity-decomposition-preservation-criterion}
A parity-transition datum in the sense of
Definition~\ref{def:transition-parity-decomposition-datum} makes each
finite transition \(T_{\nu+1,\nu,\gamma}\) an even linear map.  Hence
\[
T_{\nu+1,\nu,\gamma}
\bigl(V_{\nu+1,\gamma,\overline p}\bigr)
\subset
V_{\nu,\gamma,\overline p},\qquad p=0,1,
\]
and the even and odd retained root-space towers are defined
separately.  The same conclusion holds for the negative-root spaces
and is compatible with the finite comparison maps \(A_\beta\).  Without
the rows of
Definition~\ref{def:transition-parity-decomposition-datum}, transition
preservation of parity decompositions is not established.
\end{proposition}

\begin{proof}
The commutator identity
\(J_{\nu,\gamma}T_{\nu+1,\nu,\gamma}
=T_{\nu+1,\nu,\gamma}J_{\nu+1,\gamma}\)
is the coordinate-free statement that the transition is even.  If
\(x\in V_{\nu+1,\gamma,\overline p}\), then
\[
J_{\nu,\gamma}T_{\nu+1,\nu,\gamma}x
=T_{\nu+1,\nu,\gamma}J_{\nu+1,\gamma}x
=(-1)^pT_{\nu+1,\nu,\gamma}x.
\]
Thus \(T_{\nu+1,\nu,\gamma}x\) lies in the
\((-1)^p\)-eigenspace of \(J_{\nu,\gamma}\), namely
\(V_{\nu,\gamma,\overline p}\).  In parity-homogeneous bases this is
equivalent to the two off-diagonal block identities
\(T_{\overline 0\overline 1}=T_{\overline 1\overline 0}=0\).  The
restricted matrices are the diagonal blocks.  The Chevalley
anti-involution and the \(A_\beta\)-comparison square are additional
finite matrix identities in the datum, so the same parity preservation
holds on negative roots and after comparison with the target
presentation.  The composition row gives functoriality of the even and
odd restricted maps.
\end{proof}

The packet
\[
\texttt{certificates/hall/transition\_parity\_decomposition\_preservation}
\]
records the current state of this datum.  Its verifier returns
\(\texttt{TRANSITION\_PARITY\_DECOMPOSITION\_PRESERVATION\_OBSTRUCTION\_VERIFIED}\):
the primitive-transition input, source and target-stage parity blocks,
parity-homogeneous bases, source and target fermion-parity involutions,
ambient transition matrices, parity-commutator identities,
off-diagonal zero identities, even and odd restricted transition
matrices, negative-root parity transport rows, comparison parity
squares, and parity-transition composition rows are not yet supplied.

\begin{definition}[Primitive-space Mittag--Leffler datum]
\label{def:primitive-space-ml-datum}
Fix a cofinal sequence \(W_1\subset W_2\subset\cdots\) as in
Definition~\ref{def:cofinal-primitive-recognition-datum}.  For each
retained Gram degree \(\gamma\) with
\(\alpha(\gamma)\in W_\nu\cup(-W_\nu)\), and for each parity
\(\overline p\), let
\[
P_{\nu,\gamma,\overline p}
:=P^{\Pi,\mathrm{red}}_{X,\gamma,\overline p}
\bigm|_{\alpha^{-1}(W_\nu\cup(-W_\nu))}
\]
denote the finite primitive Hall space after normal-ordered descent and
before quotienting by the Hopf radical.  When
\(W_\nu\subset W_{\nu'}\), transition preservation of primitive
subspaces gives a restricted map
\[
t_{\nu'\nu,\gamma,\overline p}:
P_{\nu',\gamma,\overline p}\longrightarrow
P_{\nu,\gamma,\overline p}.
\]
A primitive-space Mittag--Leffler datum consists of:
\begin{enumerate}
\item primitive-space rows for every
\((\nu,\gamma,\overline p)\), including a finite basis, parity, sign,
target root, and finite rank;
\item restricted primitive transition matrices \(t_{\nu'\nu}\), together
with identity and composition rows;
\item coverage rows for every retained positive and negative root
degree, for both parities, on the cofinal window sequence;
\item for every base window \(W_{\nu_0}\), every retained
\((\gamma,\overline p)\), and every later tower containing that degree,
an image-stabilization witness: there is \(\nu_1\ge \nu_0\) such that
for all \(\nu\ge\nu_1\),
\[
\operatorname{im}\bigl(
P_{\nu,\gamma,\overline p}\to
P_{\nu_0,\gamma,\overline p}\bigr)
=
\operatorname{im}\bigl(
P_{\nu_1,\gamma,\overline p}\to
P_{\nu_0,\gamma,\overline p}\bigr),
\]
with the stabilized image written in the chosen basis of
\(P_{\nu_0,\gamma,\overline p}\);
\item the defect rows
\[
\mathrm{strict}=1,\qquad \mathrm{ML}=1,\qquad
\operatorname{rank}R^1\varprojlim_\nu
P_{\nu,\gamma,\overline p}=0
\]
for every retained primitive-space tower, and compatibility rows for
the \(A_\beta\)-comparison maps.
\end{enumerate}
Transition preservation of primitive subspaces, finite rank, equality of
signed multiplicities, target-window containment, radical-transition
data, PBW-transition data, and denominator-product identities are not a
primitive-space Mittag--Leffler datum.
\end{definition}

\begin{proposition}[Vanishing of \(R^1\varprojlim\) on primitive spaces;
\textit{proved}]
\label{prop:primitive-space-lim1-vanishing-criterion}
For every tower \(P_{\nu,\gamma,\overline p}\) carrying the datum of
Definition~\ref{def:primitive-space-ml-datum}, the inverse system of
primitive spaces is Mittag--Leffler.  Consequently
\[
R^1\varprojlim_\nu P_{\nu,\gamma,\overline p}=0
\]
for every retained Gram degree, target root, sign, and parity.  Without
the rows of Definition~\ref{def:primitive-space-ml-datum}, the passage
from finite primitive Hall spaces to the completed primitive source is
not exact.
\end{proposition}

\begin{proof}
The primitive-space rows and restricted transition matrices define an
inverse system of finite-rank vector spaces once the identity and
composition rows have zero defect.  The image-stabilization row is
precisely the Mittag--Leffler condition: for each fixed base window,
the images of all sufficiently large windows inside the base primitive
space are constant.  The standard derived inverse-limit criterion for
inverse systems of vector spaces therefore gives
\(R^1\varprojlim=0\) on each primitive-space tower.  The conclusion is
degreewise and paritywise.  It cannot be inferred from pointwise
primitive kernels, dimension counts, signed Borcherds multiplicities, or
transition preservation without the stabilized-image rows.
\end{proof}

The packet
\[
\texttt{certificates/hall/primitive\_space\_lim1\_vanishing}
\]
records the current state of this datum.  Its verifier returns
\(\texttt{PRIMITIVE\_SPACE\_LIM1\_VANISHING\_OBSTRUCTION\_VERIFIED}\):
the cofinal root-window sequence, primitive-space rows, primitive bases,
finite-rank rows, parity and sign coverage rows, restricted transition
matrices, transition-functoriality rows, image subspace rows,
image-stabilization witnesses, strict Mittag--Leffler rows, zero
\(R^1\varprojlim\) rows, degreewise coverage rows, and comparison-tower
compatibility rows are not yet supplied.

\begin{definition}[Pairing-kernel Mittag--Leffler datum]
\label{def:pairing-kernel-ml-datum}
Assume the primitive-space rows of
Definition~\ref{def:primitive-space-ml-datum} have been supplied.  For
each retained \((\nu,\gamma,\overline p)\), a finite Hopf-pairing row is
a matrix for
\[
\lambda_{\nu,\gamma,\overline p}:
P_{\nu,\gamma,\overline p}\longrightarrow
P_{\nu,-\gamma,\overline p}^{\vee},
\qquad
x\longmapsto \langle x,-\rangle_\nu ,
\]
together with its opposite transpose
\[
\rho_{\nu,\gamma,\overline p}:
P_{\nu,-\gamma,\overline p}\longrightarrow
P_{\nu,\gamma,\overline p}^{\vee},
\qquad
y\longmapsto \langle -,y\rangle_\nu .
\]
Put
\[
K^{\ell}_{\nu,\gamma,\overline p}:=\ker\lambda_{\nu,\gamma,\overline p},
\qquad
K^{r}_{\nu,\gamma,\overline p}:=\ker\rho_{\nu,\gamma,\overline p}.
\]
The Hopf radical is formed later from these pairing kernels and the
radical-intersection/quotient data; a pairing-kernel datum by itself is
not a radical quotient.

A pairing-kernel Mittag--Leffler datum consists of:
\begin{enumerate}
\item finite Hopf-pairing matrices \(\lambda_{\nu,\gamma,\overline p}\)
and \(\rho_{\nu,\gamma,\overline p}\), with Hopf adjointness,
Frobenius-cyclicity, and quotient-nondegeneracy identity rows marking
that these are the pairings used in recognition;
\item finite bases and ranks for every left kernel
\(K^\ell_{\nu,\gamma,\overline p}\) and right kernel
\(K^r_{\nu,\gamma,\overline p}\);
\item transition maps between the left kernels and between the right
kernels, induced by the primitive transition matrices and the
pairing-transition identity, with identity and composition rows;
\item coverage rows for every retained positive and negative root
degree, both parities, and both kernel sides;
\item for every base window \(W_{\nu_0}\), retained
\((\gamma,\overline p)\), and side \(s\in\{\ell,r\}\), an
image-stabilization witness: there is \(\nu_1\ge\nu_0\) such that for
all \(\nu\ge\nu_1\),
\[
\operatorname{im}\bigl(K^s_{\nu,\gamma,\overline p}\to
K^s_{\nu_0,\gamma,\overline p}\bigr)
=
\operatorname{im}\bigl(K^s_{\nu_1,\gamma,\overline p}\to
K^s_{\nu_0,\gamma,\overline p}\bigr),
\]
with the stabilized image written in a finite kernel basis;
\item the defect rows
\[
\mathrm{strict}=1,\qquad \mathrm{ML}=1,\qquad
\operatorname{rank}R^1\varprojlim_\nu
K^s_{\nu,\gamma,\overline p}=0
\]
for every retained pairing-kernel tower, and compatibility rows with
the \(A_\beta\)-comparison to the target invariant-pairing kernels.
\end{enumerate}
Hopf-pairing matrices, rank equalities, quotient non-degeneracy,
primitive-space Mittag--Leffler exactness, radical-transition
preservation, target invariant kernels, and denominator identities are
not a pairing-kernel Mittag--Leffler datum.
\end{definition}

\begin{proposition}[Vanishing of \(R^1\varprojlim\) on pairing kernels;
\textit{proved}]
\label{prop:pairing-kernel-lim1-vanishing-criterion}
For every tower \(K^s_{\nu,\gamma,\overline p}\), \(s\in\{\ell,r\}\),
carrying the datum of
Definition~\ref{def:pairing-kernel-ml-datum}, the inverse system of
pairing kernels is Mittag--Leffler.  Consequently
\[
R^1\varprojlim_\nu K^s_{\nu,\gamma,\overline p}=0
\]
for every retained Gram degree, target root, sign, parity, and kernel
side.  Without the rows of
Definition~\ref{def:pairing-kernel-ml-datum}, kernel formation for the
Hopf pairing need not commute with passage from finite windows to the
completed primitive source.
\end{proposition}

\begin{proof}
The pairing-kernel rows and restricted transition matrices define an
inverse system of finite-rank vector spaces once identity and
composition defects vanish.  The image-stabilization row is exactly the
Mittag--Leffler condition for this kernel system.  The standard derived
inverse-limit criterion for inverse systems of vector spaces gives
\(R^1\varprojlim=0\) for each left and right pairing-kernel tower.  This
is a statement about the kernels as subspaces with their restricted
transition maps.  It does not follow from the existence of a Hopf
pairing matrix, from rank equality, from non-degeneracy after quotient,
or from Mittag--Leffler exactness of the ambient primitive spaces without
strict kernel-transition and stabilized-image rows.
\end{proof}

The packet
\[
\texttt{certificates/hall/pairing\_kernel\_lim1\_vanishing}
\]
records the current state of this datum.  Its verifier returns
\(\texttt{PAIRING\_KERNEL\_LIM1\_VANISHING\_OBSTRUCTION\_VERIFIED}\):
the cofinal root-window sequence, primitive-space input, Hopf-pairing
matrix rows, Hopf-pairing identity rows, left and right kernel rows,
kernel bases, finite-rank rows, pairing-transition input, restricted
kernel transition matrices, transition-functoriality rows, image
subspace rows, image-stabilization witnesses, strict Mittag--Leffler
rows, zero \(R^1\varprojlim\) rows, degreewise coverage rows, and
comparison-kernel compatibility rows are not yet supplied.

\begin{definition}[PBW associated-graded Mittag--Leffler datum]
\label{def:pbw-associated-graded-ml-datum}
Assume the radical quotient transition data and the PBW-transition
datum of Definition~\ref{def:transition-pbw-filtration-datum} have been
supplied.  Let
\[
L_\nu:=P^{\Pi,\mathrm{red}}_{X}|_{\alpha^{-1}(S_{W_\nu})}
/\operatorname{Rad}^{\Pi}_{X,\nu}
\]
be the finite primitive radical quotient on the window \(W_\nu\).  For
each PBW degree \(d\), retained Gram degree \(\chi\), and parity
\(\overline p\), put
\[
G^{\mathrm{PBW}}_{\nu,d,\chi,\overline p}
:=\Bigl(\operatorname{gr}^{\mathrm{PBW}}_dU(L_\nu)\Bigr)_{\chi,\overline p}.
\]
The associated-graded transition maps of
Proposition~\ref{prop:transition-pbw-filtration-preservation-criterion}
give maps
\[
g_{\nu'\nu,d,\chi,\overline p}:
G^{\mathrm{PBW}}_{\nu',d,\chi,\overline p}
\longrightarrow
G^{\mathrm{PBW}}_{\nu,d,\chi,\overline p}
\qquad (\nu\le\nu').
\]
A PBW associated-graded Mittag--Leffler datum consists of:
\begin{enumerate}
\item finite PBW filtration rows \(F^{\mathrm{PBW}}_{\le d}U(L_\nu)\),
lower filtration rows \(F^{\mathrm{PBW}}_{\le d-1}U(L_\nu)\), and
quotient rows for every
\(G^{\mathrm{PBW}}_{\nu,d,\chi,\overline p}\);
\item ordered word bases, relation-rewrite rows, associated-graded
bases, and associated-graded rank rows;
\item associated-graded transition matrices \(g_{\nu'\nu}\), together
with PBW-symbol intertwining, \(A_\beta\)-PBW comparison squares,
identity rows, and composition rows;
\item coverage rows for every retained PBW degree, Gram degree, parity,
side, and cofinal window;
\item for every base window \(W_{\nu_0}\), every retained
\((d,\chi,\overline p)\), and every side, an image-stabilization
witness: there is \(\nu_1\ge\nu_0\) such that for all \(\nu\ge\nu_1\),
\[
\operatorname{im}\bigl(
G^{\mathrm{PBW}}_{\nu,d,\chi,\overline p}\to
G^{\mathrm{PBW}}_{\nu_0,d,\chi,\overline p}\bigr)
=
\operatorname{im}\bigl(
G^{\mathrm{PBW}}_{\nu_1,d,\chi,\overline p}\to
G^{\mathrm{PBW}}_{\nu_0,d,\chi,\overline p}\bigr),
\]
with the stabilized image written in a finite associated-graded basis;
\item the defect rows
\[
\mathrm{strict}=1,\qquad \mathrm{ML}=1,\qquad
\operatorname{rank}R^1\varprojlim_\nu
G^{\mathrm{PBW}}_{\nu,d,\chi,\overline p}=0
\]
for every retained PBW associated-graded tower.
\end{enumerate}
PBW filteredness, associated-graded rank equality, target PBW rows,
Hilbert series, formal PBW pushforward, denominator identities, and
primitive or pairing-kernel Mittag--Leffler exactness are not a PBW
associated-graded Mittag--Leffler datum.
\end{definition}

\begin{proposition}[Vanishing of \(R^1\varprojlim\) on PBW associated
gradeds; \textit{proved}]
\label{prop:pbw-associated-graded-lim1-vanishing-criterion}
For every tower
\(G^{\mathrm{PBW}}_{\nu,d,\chi,\overline p}\) carrying the datum of
Definition~\ref{def:pbw-associated-graded-ml-datum}, the inverse system
of PBW associated-graded pieces is Mittag--Leffler.  Consequently
\[
R^1\varprojlim_\nu
G^{\mathrm{PBW}}_{\nu,d,\chi,\overline p}=0
\]
for every retained PBW degree, Gram degree, parity, and side.  Without
the rows of Definition~\ref{def:pbw-associated-graded-ml-datum}, PBW
associated gradeds need not commute with passage from finite windows to
the completed enveloping algebra.
\end{proposition}

\begin{proof}
The associated-graded piece rows and associated-graded transition
matrices define an inverse system of finite-rank vector spaces once the
identity and composition rows have zero defect.  The
image-stabilization row is precisely the Mittag--Leffler condition on
that inverse system.  The standard derived inverse-limit criterion for
inverse systems of vector spaces gives \(R^1\varprojlim=0\) on every
retained PBW associated-graded tower.  Filteredness alone gives maps of
filtered pieces; it does not imply that the images on the quotients
\(F_{\le d}/F_{\le d-1}\) stabilize.  Rank equality and Hilbert-series
agreement likewise do not compute the derived inverse limit.
\end{proof}

The packet
\[
\texttt{certificates/hall/pbw\_associated\_graded\_lim1\_vanishing}
\]
records the current state of this datum.  Its verifier returns
\(\texttt{PBW\_ASSOCIATED\_GRADED\_LIM1\_VANISHING\_OBSTRUCTION\_VERIFIED}\):
the cofinal root-window sequence, radical-quotient input,
PBW-transition input, PBW filtration rows, ordered word bases,
relation-rewrite rows, associated-graded piece rows, graded bases,
associated-graded rank rows, associated-graded transition maps,
PBW-symbol intertwining rows, \(A_\beta\)-PBW squares,
transition-functoriality rows, image subspace rows,
image-stabilization witnesses, strict Mittag--Leffler rows, zero
\(R^1\varprojlim\) rows, degreewise coverage rows, and
comparison-graded compatibility rows are not yet supplied.

\begin{proposition}[Strict Mittag--Leffler completion for primitive
recognition; \textit{proved}]
\label{prop:strict-ml-completion-primitive-recognition}
Let
\(\{W_\nu\}_{\nu\ge 1}\) be a cofinal increasing sequence of
downward-saturated finite root windows.  For each \(\nu\), let
\[
A_\nu:\ P^{\Pi,\mathrm{red}}_{X}|_{\alpha^{-1}(S_{W_\nu})}
\xrightarrow{\sim}
\mathfrak g_{\Delta_5}|_{S_{W_\nu}}
\]
be the finite-window isomorphism of
Proposition~\ref{prop:finite-window-presentation-recognition}.  Suppose
the transition maps for \(W_{\nu+1}\to W_\nu\) satisfy the following
strict Mittag--Leffler conditions:
\begin{enumerate}
\item the source primitive-space towers of
Definition~\ref{def:primitive-space-ml-datum} and the target root-space
parity towers are Mittag--Leffler, and \(A_\nu\) commutes with the
transition maps on each parity piece;
\item the pairing-kernel towers of
Definition~\ref{def:pairing-kernel-ml-datum}, the kernel towers of the
finite presentation maps, the Hopf radical intersection towers, and the
relation-radical quotient towers are strict subsystems: transition maps
carry each subspace onto the corresponding image and induce the quotient
transition maps;
\item the PBW filtrations are strict under transition, and the induced
towers on the associated graded pieces carry the
Mittag--Leffler datum of
Definition~\ref{def:pbw-associated-graded-ml-datum};
\item the triangular finite-window towers for positive-negative
cancellation have \(R^1\varprojlim=0\).
\end{enumerate}
Then inverse limit is exact on the finite recognition diagram.  In
particular,
\[
\varprojlim_\nu
\bigl(P^{\Pi,\mathrm{red}}_{X}|_{\alpha^{-1}(S_{W_\nu})}
/\operatorname{Rad}^{\Pi}_{X,\nu}\bigr)
\cong
\bigl(\varprojlim_\nu
P^{\Pi,\mathrm{red}}_{X}|_{\alpha^{-1}(S_{W_\nu})}\bigr)
\big/
\overline{\varprojlim_\nu\operatorname{Rad}^{\Pi}_{X,\nu}},
\]
the completed kernel is the closure of the finite kernels, the
completed PBW associated graded is the inverse limit of the finite PBW
associated gradeds, and the maps \(A_\nu\) assemble to an isomorphism
of completed Gram-graded Lie superalgebras
\[
\varprojlim_\nu
P^{\Pi,\mathrm{red}}_{X}|_{\alpha^{-1}(S_{W_\nu})}
\xrightarrow{\sim}
\varprojlim_\nu \mathfrak g_{\Delta_5}|_{S_{W_\nu}}.
\]
Without these strict Mittag--Leffler conditions, finite-window
isomorphisms do not by themselves identify the completed primitive
source with the completed BKM target.
\end{proposition}

\begin{proof}
For inverse systems of vector spaces, the derived exact sequence for
\(\varprojlim\) shows that exactness of
\[
0\to K_\nu\to V_\nu\to Q_\nu\to0
\]
passes to inverse limits exactly when \(R^1\varprojlim K_\nu=0\), or
more generally when the tower is Mittag--Leffler.  Apply this to the
presentation-kernel tower, the Hopf-radical tower, the quotient tower,
and each even/odd root-space tower.  Strictness says that taking
subobjects, quotients, and finite PBW filtered pieces commutes with the
transition maps; hence the inverse-limit quotient is the quotient by
the closure of the finite radicals, and the inverse-limit kernel is the
closure of the finite kernels.  The strict filtered statement gives
\[
\operatorname{gr}_{\mathrm{PBW}}\varprojlim_\nu U(P_\nu)
\cong
\varprojlim_\nu\operatorname{gr}_{\mathrm{PBW}}U(P_\nu)
\]
on every retained degree.

The maps \(A_\nu\) commute with the transition maps, so they form an
isomorphism of inverse systems.  Since inverse limit is a functor and
the kernel and cokernel towers of \(A_\nu\) are zero, the inverse-limit
map has zero kernel and cokernel.  If the Mittag--Leffler hypotheses are
removed, the \(R^1\varprojlim\) term can contribute a hidden kernel,
cokernel, radical element, PBW class, or parity-cancelling summand
which is invisible on each fixed finite window.  This is exactly the
completion failure prohibited by (R8).
\end{proof}

\begin{definition}[Finite source proof tables]
\label{def:finite-source-proof-tables}
After choosing homogeneous bases on a finite window \(W_\nu\), the
finite-window primitive-recognition datum of
Definition~\ref{def:cofinal-primitive-recognition-datum} is equivalent
to the following source proof tables:
\[
\begin{gathered}
\mathrm{degrees},\ \mathrm{parity\_blocks},\
\mathrm{basis\_provenance},\
\mathrm{simple\_representatives},\
M,\ D,\ \mathrm{unit\_counit},\ B,\ G,\
\mathrm{hopf\_pairing\_identities},\ K,\ Q,\ A,\\
\mathrm{hall\_bialgebra\_identities},\
\mathrm{radical\_ideal\_coideal},\
\mathrm{relation\_rows},\
\mathrm{no\_extra},\
\mathrm{generation},\
\mathrm{pbw},\
\mathrm{transitions},\
\mathrm{A}_{\beta}\mathrm{\_comparison\_maps}.
\end{gathered}
\]
When the same packet is used to assert the current-envelope statement
\(\Omega_E^{\mathrm{ch}}C_X\simeq\mathsf{Den}_{\Delta_5,E}\), it must
also contain the finite Koszul comparison tables
\[
\mathrm{koszul\_cones},\qquad
\mathrm{koszul\_comparison\_identities},\qquad
\mathrm{koszul\_transition\_ml}
\]
of Definition~\ref{def:finite-koszul-comparison-datum}.  They are not
implied by the \(A_\beta\)-comparison maps: \(A_\beta\) records the
primitive radical-quotient comparison, whereas the Koszul tables record
the acyclic cones, compatibility homotopies, and strict
Mittag--Leffler systems of the chiral cobar comparison.

Their mathematical meanings are fixed as follows.  The degree table
records the Gram labels \(\gamma\), their signs, window membership, and
compact source-admissibility.  The provenance table records the retained
charge lift, source stratum, vanishing-cycle or intersection-complex
summand, reduced orientation, quotient orientation,
Thom--Sebastiani transport, finite-stabilizer linearization, protected
integration, and transition compatibility of every basis vector.  The
parity table records the two ranks in each parity, not only their
difference.  The simple-representative table names the source basis
vectors realizing \(e_i^X,f_i^X,h_i^X\), together with their parity and
target reference.  The matrices \(M,D,B\) record, respectively, the
source Hall product, source coproduct, and supercommutator bracket on
the chosen source bases.  The \(\mathrm{unit\_counit}\) table records
the unit and counit maps.  The Hall bialgebra identity table records the
vanishing of the finite defects for associativity, coassociativity,
left and right unit, left and right counit, product-coproduct
compatibility, and primitive closure of the supercommutator.  The
matrix \(G\) records the Hopf pairing.  The Hopf-pairing identity table
records the finite defects for Hopf adjointness
\(\langle m(x,y),z\rangle=\langle x\otimes y,\Delta z\rangle\), the
Frobenius cyclic identity for the primitive bracket, and
non-degeneracy of the induced pairing on the radical quotient.  The
matrices \(K,Q\) record the radical kernel and a quotient splitting.
The matrix \(A\) records the comparison from the source radical quotient
to the imported target basis; target labels are allowed only in this
codomain.

Each row is \emph{geometrically generated}.  It carries a source
identifier for the compact correspondence, reduced stratum, vanishing-
cycle or intersection-complex summand, orientation transport,
Hopf-pairing construction, radical computation, quotient splitting,
comparison map, or transition morphism from which the row is obtained,
together with a proof reference for the relevant theorem, equation,
finite computation, or source certificate.  A row whose matrices satisfy
the displayed rank identities but are not obtained from these
compact-source data is not a row of the finite source proof tables.  In
particular, arbitrary matrices with the correct dimensions, target-only
rows, signed-only rows, and status-only rows do not define the compact
source primitive-recognition datum.

The remaining tables record identities.  The
\(\mathrm{radical\_ideal\_coideal}\) rows state that \(K\) is stable as
a Lie ideal and as a coproduct coideal.  The
\(\mathrm{relation\_rows}\) table records the Cartan, Chevalley, real
Serre, Borcherds orthogonality, and super sign relations.  The
\(\mathrm{no\_extra}\) table records the finite rank identity
\(\ker\pi_\nu=(\overline{J_{\mathrm{BK}}+\operatorname{Rad}_{\mathrm{GN}}})_{\le W_\nu}\).
The \(\mathrm{generation}\) table records, for each source basis vector
in the radical quotient, a bracket-word span by the Cartan and simple
primitive representatives with all intermediate degrees remaining in
\(W_\nu\).  The \(\mathrm{pbw}\) table records the associated-graded PBW
rank comparison.  The \(\mathrm{transitions}\) table records strictness
and Mittag--Leffler exactness under \(W_{\nu+1}\to W_\nu\).  The
\(\mathrm{A}_{\beta}\)-comparison table records that the comparison
matrix \(A\) intertwines the source bracket, coproduct, pairing, radical
quotient, and PBW filtration with the target blocks.
Admissibility of these finite tables is numerical: every defect-rank
row above must have rank \(0\), and every rank-comparison row must have
equal ranks on the two sides and on their combined span.  A non-zero
defect rank or a failed rank equality is a failed finite-window
recognition datum.  The identity rows are also exhaustive.  In the
finite source proof tables the canonical labels are
\[
\begin{gathered}
\texttt{unit\_left},\texttt{unit\_right},\texttt{counit\_left},
\texttt{counit\_right},\texttt{associativity},\texttt{coassociativity},
\texttt{bialgebra\_compatibility},\texttt{primitive\_closure},\\
\texttt{hopf\_adjointness},\texttt{frobenius\_cyclic},
\texttt{quotient\_nondegenerate},\quad
\texttt{lie\_ideal},\texttt{coproduct\_coideal},\\
\texttt{cartan},\texttt{chevalley},\texttt{real\_serre},
\texttt{borcherds\_orthogonality},\texttt{super\_sign},\quad
\texttt{bracket},\texttt{coproduct},\texttt{pairing},
\texttt{radical\_quotient},\texttt{pbw}.
\end{gathered}
\]
Omitting one family is a failed finite-window recognition datum.
If the current-envelope comparison is also asserted, the Koszul table
families listed above are part of the same finite packet; omitting one
of them is a failed finite Koszul comparison datum, even if
\((\mathrm{R0})\)--\((\mathrm{R8})\) hold algebraically.
\end{definition}

For the compact \(K3\times E\) Hall source packet used in this
manuscript, the companion obstruction ledger
\[
\texttt{certificates/sources/k3e\_compact\_hall/blocked\_obligations.csv}
\]
records precisely the missing source artifacts.  Its verifier returns
\(\texttt{COMPACT\_HALL\_OBSTRUCTION\_LEDGER\_VERIFIED}\), while still
recording \(\texttt{compact\_source\_recognition=false}\) and
\(\texttt{mathematical\_certification=false}\).  The ledger is therefore
not evidence for primitive recognition.  It is the finite list of
objects which must be supplied before \(P_X^{\Pi,\mathrm{red}}\) exists
as a constructed compact-source primitive Hall Lie superalgebra:
the coproduct \(D\), unit and counit, bialgebra compatibility, primitive
closure, Hopf pairing, radical ideal/coideal, no-extra kernel equality,
PBW comparison, strict Mittag--Leffler transition, the
\(A_\beta\)-intertwining rows, and the finite Koszul comparison rows.

\begin{proposition}[The finite source proof tables are exactly the
finite recognition datum; \textit{proved}]
\label{prop:finite-source-proof-tables-equivalence}
For a fixed downward-saturated finite window \(W_\nu\), a populated
system of finite source proof tables in the sense of
Definition~\ref{def:finite-source-proof-tables}, with all stated rank
and intertwining identities verified over the chosen coefficient ring
and with every row geometrically generated from the compact-source
correspondences and orientation/Hopf data named there,
is equivalent to clauses \textup{(R0)}--\textup{(R8)} of
Definition~\ref{def:cofinal-primitive-recognition-datum} after a choice
of homogeneous bases.  In particular, a header-only table system, or a
table system whose rows contain status text but no mathematical payload,
or a table system of arbitrary matrices without compact-source
provenance, does not define a finite-window primitive-recognition datum.
\end{proposition}

\begin{proof}
Choosing homogeneous bases turns each finite-dimensional source
primitive space, pairing block, radical, quotient, and comparison map
into matrices, but only after the corresponding finite Hall
correspondence, reduced orientation, Hopf pairing, radical quotient, and
transition maps have been constructed.  The row-level geometric source
and proof reference record this construction before coordinates are
chosen.  The degree and provenance tables are precisely (R0).
The simple-representative table is (R1).  The parity table is (R3).
The representative rows together with the relation table are
(R1)--(R2).
The \(M,D,\eta,\epsilon,B\) matrices and the Hall bialgebra identity
table are the product-coproduct and primitive-closure part of (R4).
The \(G\) matrix and the Hopf-pairing identity table are the pairing
adjointness, Frobenius, and quotient non-degeneracy part of (R4).  The
\(K,Q\) matrices and the
\(\mathrm{radical\_ideal\_coideal}\) identities are the radical part of
(R4).  The
\(\mathrm{no\_extra}\) rank identity is (R5).  The generation span table
is (R6).  The PBW table is (R7), and the transition table is (R8).  The
\(\mathrm{A}_{\beta}\)-comparison rows assert that the source quotient
is compared to the same target blocks in every operation, so the matrix
statements are independent of the chosen coordinates.  Conversely, any
finite-window datum satisfying (R0)--(R8) yields these tables after
choosing bases and writing the finite maps and identities in matrix
form, with the geometric source and proof reference inherited from the
finite maps themselves.  Empty rows have no source vectors, maps, ranks,
or identities; arbitrary matrices with no compact-source origin have no
Hall correspondence, orientation datum, pairing construction, or
transition map.  They satisfy none of the clauses.
\end{proof}

\begin{proposition}[Global recognition is exactly the cofinal datum;
\textit{proved}]
\label{prop:global-recognition-cofinal-datum}
Assume (S1)--(S4) of Theorem~\ref{thm:primitive-recognition}.  Then
the global source-side assertions (S5)--(S10) are equivalent to
a cofinal finite-window primitive-recognition datum
$\{W_\nu,(\textup{R0})\text{--}(\textup{R8})\}$ in the sense of
Definition~\ref{def:cofinal-primitive-recognition-datum}, and
conversely.
\end{proposition}

\begin{proof}
The two sides record the same data: each $W_\nu$ is a saturated
finite window in (S5)--(S9), and (R3) is exactly the
two-integer parity equality of (S10) on $W_\nu$.  The Hopf pairing
matrices, kernel equalities, PBW filtrations, no-extra-relations
relations, and transition strictness are exactly the cofinal
contents of (S5)--(S9), and the cusp-positive arithmetic side of
$f(nm,l)$ is recovered from the parity equalities (R3) on the active
target degrees.  Conversely, given (R0)--(R8) on a cofinal system,
the strict Mittag--Leffler completion theorem
Proposition~\ref{prop:strict-ml-completion-primitive-recognition}
gives the global statements (S5)--(S9), and (R3) on a cofinal
exhaustion gives (S10).
\end{proof}

\subsubsection*{Recognition criterion at the window $W_{\le 3}$}

For a finite window $W$ with $S_W=W\cup(-W)\cup\{0\}$, a compact
source datum consists of compact-source bases
$b^X_{R,\gamma_\beta,\overline p,a}$, parity pieces, product,
coproduct, bracket, pairing, radical, and quotient matrices
$M,D,B,G,K,Q$, comparison maps $A_{\beta,\overline p}$, relation
identities, the no-extra kernel equality, PBW associated gradeds,
and strict Mittag--Leffler transition maps.  The target labels
$e_i,E_{ij},u_{ij,r},T_i,M_{ij,r},w_s$ occur only in target blocks
or as codomain coordinates of $A_{\beta,\overline p}$.

The window $W_{\le 3}$ is the first arithmetic window, not a
relation-closed degree window.  A relation-closed degree test must include
real-Serre terminal rows, doubled isotropic rows, $\delta_{123}$-real
strings, odd--odd $\delta_{123}$ rows, and complementary
height-seven strings, or pass to the downward-saturated closure.  In
the target-reference packet \(\texttt{wle3\_target\_parity}\),
\(W_{\le3}\) has the verified rows \(\delta_i:1|0\),
\(a_{ij}:10|0\), \(2\delta_i+\delta_j:1|0\), and
\(\delta_{123}:29|93\); its scalar firewall excludes source parity,
Hall brackets, pairing radicals, PBW comparison, and primitive
recognition.
The finite packet \(\texttt{parity\_pushforward}\) verifies a different
formal statement: once finite source parity blocks are supplied, the
additive normal-ordered Gram pushforward preserves the even and odd
blocks separately.  It supplies neither the source parity dimensions
required in (R3) nor the Gritsenko--Nikulin/Kac target parity equality;
those remain primitive-recognition data.
The finite packet \(\texttt{pbw\_pushforward}\) is likewise formal: once
a finite PBW filtration is supplied, additive normal-ordered Gram
pushforward preserves its filtered pieces and associated graded.  It
does not supply the compact-source PBW comparison required in (R7), the
target PBW rows, or the no-extra-relations theorem.
The finite packet \(\texttt{triangular\_pushforward}\) is the analogous
formal statement for a supplied positive/Cartan/negative splitting: the
normal-ordered pushforward preserves the three direct-sum parts and the
displayed bracket signs.  It does not supply the exact triangular
completion required in (R8), target triangular completion, or the
Mittag--Leffler exactness of triangular finite-window systems.
The finite target degree-closure packet
\(\texttt{wrel\_degree\_closure}\) verifies that the enlarged target
degree set has exactly the \(35\) rows
\[
W_{\le3}^{act}\cup R_4\cup I_4\cup C_{\le7}\cup\{2\delta_{123}\},
\]
with \(R_4\) and \(C_{k,5}\) terminal zero rows.  Its
downward-saturation audit finds exactly three nonzero proper
subdegree defects, namely
\[
D_k=\delta_k+2a_{ij},\qquad \{i,j,k\}=\{1,2,3\},
\]
below \(2\delta_{123}\), each with signed multiplicity \(-513\).
This is still only the target degree closure.
In the larger target degree window,
Remark~\ref{rem:bkm-relation-window-target-parity}
supplies full parity only for \(2a_{ij}\), \(C_{k,3}\), \(C_{k,4}\),
and \(C_{k,5}\).  The rows \(C_{k,2}\), \(D_k\), and
\(2\delta_{123}\) remain signed target rows; they cannot serve as
source comparison codomains until the finite target presentation is
computed in those degrees.  The base-string terminal degree set is
therefore not yet a primitive-recognition window in the sense of
(R0)--(R8).
The A071 target-presentation verifier records this boundary: it checks
the \(2a_{ij}\), \(C_{k,3}\), \(C_{k,4}\), and \(C_{k,5}\) parity rows,
and rejects \(C_{k,2}\), \(D_k\), and \(2\delta_{123}\) as target
basis, pairing, radical, or PBW rows.  The failure is structural, not
a missing scalar coefficient.  The minimal nonzero downward-saturation
extension \(W_{\mathrm{rel}}^+\cup\{D_1,D_2,D_3\}\) is now
target-recorded; a primitive-recognition window must still supply the
missing target-presentation and compact-source data.  The
post-saturation real-string audit records \(21\) terminal codomains
outside the saturated extension, \(10\) with nonzero signed
multiplicity.  The terminal-layer audit groups these codomains into
\(21\) distinct terminal degrees; adjoining them creates \(114\)
nonzero subdegree-defect incidences across \(26\) distinct subdegrees.
The next target-degree computation adjoins exactly these \(26\) rows and
gives an \(85\)-row degree set with zero remaining nonzero
proper-subdegree defect.  This closes downward saturation for the
displayed terminal layer, but full finite relation closure remains a
further target closure problem: if these \(26\) rows are promoted to
target-generator rows, the real-root strings force \(55\) terminal
checks, \(54\) terminal codomains outside the \(85\)-row set, and
\(12\) missing codomains with nonzero signed multiplicity.  Adjoining
those \(55\) terminal codomains would create \(546\) nonzero
proper-subdegree defect incidences across \(98\) distinct subdegrees.
Adjoining those \(98\) rows gives a \(237\)-row target degree set with
zero remaining nonzero proper-subdegree defect for that layer.  If this
\(98\)-row layer is promoted to target-generator rows, the real-root
strings force \(206\) terminal checks, \(182\) new terminal degrees, and
\(24\) missing nonzero terminal degrees.  The exact terminal-layer table
records all \(206\) codomains, and the exact saturation-extension table
records the \(50\) distinct nonzero proper subdegrees created by
adjoining them, with \(624\) incidences.  Adjoining both layers gives a
\(469\)-row target degree set with zero remaining nonzero
proper-subdegree defects for that layer.  Promoting the newly adjoined
\(50\)-row layer forces \(96\) real-string terminal checks, \(90\) new
terminal degrees, no missing nonzero terminal degree, and a \(12\)-row
nonzero saturation extension with \(44\) incidences.  Adjoining these
rows gives a \(571\)-row target degree set with zero remaining nonzero
proper-subdegree defects for that layer.  Promoting the newly adjoined
\(12\)-row layer forces \(20\) real-string terminal checks, all with
zero signed terminal multiplicity, and a \(6\)-row nonzero saturation
extension with \(8\) incidences.  Adjoining these rows gives a
\(597\)-row target degree set with zero remaining nonzero
proper-subdegree defects for that layer.  Promoting the newly adjoined
\(6\)-row layer forces \(10\) real-string terminal checks, all with
zero signed terminal multiplicity, and a \(6\)-row nonzero saturation
extension with \(18\) incidences.  Adjoining these rows gives a
\(613\)-row target degree set with zero remaining nonzero
proper-subdegree defects for that layer.  Promoting the next newly
adjoined \(6\)-row layer forces \(10\) real-string terminal checks, all
with zero signed terminal multiplicity, and a \(6\)-row nonzero
saturation extension with \(8\) incidences.  Adjoining these rows gives
a \(629\)-row target degree set with zero remaining nonzero
proper-subdegree defects for that layer.  Promoting the next newly
adjoined \(6\)-row layer forces \(10\) real-string terminal checks, all
with zero signed terminal multiplicity, and a \(6\)-row nonzero
saturation extension with \(18\) incidences.  Adjoining these rows gives
a \(645\)-row target degree set with zero remaining nonzero
proper-subdegree defects for that layer.  The tail is not a finite
accident.  For every even \(N\ge 14\), put
\[
\begin{aligned}
A_N&=\{(0,N,N-1),(0,N,N),(0,N+1,N),
      (N,0,N-1),(N,0,N),(N+1,0,N)\},\\
B_N&=\{(0,N,N+1),(0,N+1,N+1),(0,N+1,N+2),
      (N,0,N+1),(N+1,0,N+1),(N+1,0,N+2)\}.
\end{aligned}
\]
The verifier
\(\texttt{wrel\_tail\_staircase}\) proves symbolically, using the
Lorentzian pairing and the \(q^0\)-term \(y^{-1}+10+y\) of
\(\phi_{0,1}\), that the real-root terminal audit and downward
saturation force
\[
A_N\Longrightarrow B_N,\qquad B_N\Longrightarrow A_{N+2}.
\]
Its finite-anchor table checks the six rows of \(A_{14}\) inside the
\(\texttt{wrel\_degree\_closure}\) packet.  Hence the target-degree
audit contains an infinite six-row staircase.  Its
\(\texttt{unbounded\_tail.csv}\) table records the subsequences
\(A_{14+2k}\) and \(B_{14+2k}\) with a strictly increasing coordinate.
The source comparison can therefore be formulated only through finite
presentation data, not by closing a finite target-degree enumeration.
To promote this target audit to the first relation-closed compact-source
theorem, one must populate the first-window packet
\[
\mathrm{window\_closure},\quad
\mathrm{target\_source\_representatives},\quad
\mathrm{chevalley\_serre\_matrices},\quad
\mathrm{hall\_boundary\_complex},\quad
\mathrm{spectral\_sequence},\quad
\mathrm{kernel\_matrices},\quad
\mathrm{pbw\_graded},\quad
\mathrm{first\_window\_theorems},\quad
\mathrm{transitions},\quad
\mathrm{scalar\_firewall}.
\]
The theorem rows include the first-window instances of \(O1\),
\(O1^+\), \(O2\), \(P_{\mathrm{fin}}\), Hall product, Hall coproduct,
Hopf pairing, Koszul comparison, PBW comparison, primitive
recognition, and Pfaffian equality.  A target relation-closed degree
set without these compact-source rows is not a source theorem.
At present the packet
\(\texttt{certificates/first\_window/k3e\_relation\_closed\_window}\)
verifies only the obstruction ledger
\(\texttt{FIRST\_WINDOW\_OBSTRUCTION\_LEDGER\_VERIFIED}\), with
\(\texttt{first\_window\_certification=false}\).  Thus the first
relation-closed source theorem remains open until these finite rows are
supplied.

\begin{proposition}[Target-admissibility obstruction for the
height-seven closure; \textit{proved}]
\label{prop:target-admissibility-obstruction-wrel}
Let \(W\) be a downward-saturated finite target-root window containing
the height-\(7\) closure \(W_{\mathrm{rel}}^+\) of
Chapter~\ref{ch:pfaffian-dirac}.  If \(W\) is a
primitive-recognition window satisfying (R0)--(R8), then the target
restriction must contain a finite target-presentation datum.  In
particular, for every row \(C_{k,2}\), \(D_k=\delta_k+2a_{ij}\), and
\(2\delta_{123}\), it must supply the full even/odd parity dimensions,
the positive-negative pairing block, the target radical block, and the
PBW associated-graded row.  The signed target identities
\[
 \sdim(C_{k,2})=108,\quad m(C_{k,2})=90,\qquad
 \sdim(D_k)=-513,\quad m(D_k)=54,\qquad
 \sdim(2\delta_{123})=4016,\quad m(2\delta_{123})=-540
\]
do not supply these data.  Hence the presently available target data
do not make the saturated extension of \(W_{\mathrm{rel}}^+\) a
primitive-recognition window.
\end{proposition}

\begin{proof}
The rows \(C_{k,2}\) and \(2\delta_{123}\) lie in
\(W_{\mathrm{rel}}^+\), and \(D_k\) lies in every downward-saturated
window containing \(W_{\mathrm{rel}}^+\), since
\[
0<D_k=\delta_k+2a_{ij}<2\delta_{123}.
\]
By
Proposition~\ref{prop:target-presentation-prerequisite}, they must be
target-presentation rows.  Clause (R3) requires the two integers
\(\rootmult_{\overline 0}^{\mathrm{GN}}(\beta)\) and
\(\rootmult_{\overline 1}^{\mathrm{GN}}(\beta)\) for every
\(\beta\in W\), and also in the negative degree by the
Borcherds--Kac Chevalley anti-involution.  A signed multiplicity gives
only their difference.  The integer \(m(\beta)\) records the
Gritsenko--Nikulin imaginary-simple contribution in the denominator
presentation; it is not the full quotient root-space parity split and
does not determine an invariant positive-negative pairing.

Clauses (R4), (R5), and (R7) then require the target pairing kernel,
radical block, and PBW associated-graded comparison in the same
degrees.  Without those target blocks the source-to-target matrices
\(A_{\beta,\overline p}\) have no defined codomain in these rows, and
the equations expressing Hopf radical equality, no-extra-relations, and
PBW compatibility are not well-formed.  Thus signed data in
\(C_{k,2}\), \(D_k\), and \(2\delta_{123}\) are not admissible target
data for (R0)--(R8).
\end{proof}

\subsubsection*{What signed dimensions miss}

\begin{lemma}[Signed dimensions are not recognition;
\textit{proved}]
\label{lem:signed-dimensions-miss}
Let $V_\bullet$ be a $\widehat\Gamma_X$-graded super vector space
with zero bracket and zero coproduct, with the same signed
superdimensions $\sdim V_\gamma=f(nm,l)$ as
$P_X^{\Pi,\mathrm{red}}$.  Then $V_\bullet$ has the same Borcherds
denominator $64^{-1}\Delta_5(2Z)$ as $\mathfrak g_{\Delta_5}$ but is
not isomorphic to it as a Lie superalgebra.
\end{lemma}

\begin{proof}
The denominator depends only on signed superdimensions; the bracket
data, no-extra-relations theorem, Hopf-radical coideal stability, and
PBW filtration are independent.  If the Cartan is mapped nontrivially,
the Chevalley relation \([e_i,f_i]=h_i\) fails in \(V_\bullet\).  If
the Cartan is killed, the root-degree Cartan identification fails.  In
either case the kernel of the free presentation map contains all
non-Cartan brackets, hence is strictly larger than
\(\overline{J_{\mathrm{BK}}+\operatorname{Rad}_{\mathrm{GN}}}\) in
degrees where the BKM bracket is nonzero.  With the zero Hopf pairing
the radical is all of \(V_\bullet\), so the radical quotient is zero.
\end{proof}

This lemma is the counterexample that prohibits the determinant route
to recognition.  It is also the structural reason the ten clauses
(S1)--(S10) do not collapse to their signed-dimension shadow:
the signed dimensions are necessary and not sufficient; (S1)--(S9)
supply the source presentation, and (S10) supplies the missing parity
split.

\begin{lemma}[Finite-window scalar exclusions;
\textit{proved}]
\label{lem:first-window-scalar-exclusions}
Let \(W\) be a finite target-root window.  None of the following data
implies a finite-window primitive-recognition datum on \(W\) in the
sense of Definition~\ref{def:cofinal-primitive-recognition-datum}:
target labels without compact representatives, signed dimensions, the
scalar trace \(Z_{BPS}^{K3\times E}\), the Pfaffian product, the
Borcherds denominator product, target PBW rows without source PBW
strictness, arbitrary matrices with the right ranks but no compact
provenance, or rows whose only content is a verification status.
\end{lemma}

\begin{proof}
Clauses (R0)--(R8) require source objects before comparison: compact
basis vectors, source representatives, Hall product and coproduct,
unit and counit, Hopf pairing, radical and quotient matrices, relation
rows, no-extra kernel equality, generation spans, PBW strictness, and
Mittag--Leffler transition maps.  Target labels give only codomain
coordinates for \(A_{\beta,\bar p}\), not the source vectors mapped to
them.  Signed dimensions and the denominator product give one integer
per degree; (R3) requires the two parity dimensions, while (R4)--(R8)
require operations and kernels.  Lemma~\ref{lem:signed-dimensions-miss}
exhibits a zero-bracket object with the same denominator and no
recognition isomorphism.  The scalar trace and the Pfaffian product
live after taking a protected scalar shadow; they forget the bracket,
coproduct, Hopf radical, PBW filtration, and finite transition system.
Target PBW rows test only the imported target presentation and do not
give source PBW strictness or a source-to-target intertwining matrix.
Arbitrary matrices with correct ranks do not supply a Hall
correspondence, reduced orientation, compact cycle, pairing
construction, or transition morphism, so they fail the geometric
provenance part of (R0) and the construction clauses of (R4)--(R8).
Status-only rows contain no vector space, map, matrix, kernel, rank, or
identity.  Hence each listed substitute omits at least one structural
component of (R0)--(R8), and none defines a finite-window
primitive-recognition datum.
\end{proof}

\subsubsection*{Obstruction classes}

The recognition lane carries the residuals
\[
\mathfrak O_{\mathrm{Rec},\nu}=
\bigl(
\mathfrak o^{(\textup{R1})}_\nu,
\ldots,
\mathfrak o^{(\textup{R8})}_\nu,
\mathfrak o^{\mathrm{tr}}_{\nu+1,\nu}
\bigr),
\]
required to vanish on every cofinal window.  The first arithmetic
window $W_{\le 3}$ is achievable from the GN target side; the first
relation-closed degree window includes the real-Serre terminal rows and the
$\delta_{123}$ odd--odd rows, but some added target rows are still only
signed target rows by
Remark~\ref{rem:bkm-relation-window-target-parity}.  The
finite-stage components of the
(R0)--(R8) criterion, including source-radical/Hopf-coideal stability,
remain source-recognition data.  Appendix~\ref{app:obstruction-discharges}
supplies the orientation lane; it does not supply the source matrices
\(M,D,B,G,K,Q\), the comparison maps \(A_{\beta,\bar p}\), the kernel
equality, the PBW associated-graded comparison, or the finite Koszul
comparison residual
\(\mathfrak O_{\mathrm{Kcmp},R}\).  Construction of the
source-internal recognition criterion at any one relation-closed degree
window beyond \(W_{\le 3}\), together with the Koszul compatibility
checks of Definition~\ref{def:finite-koszul-comparison-datum}, is the
unresolved compact-source datum of Chapter~\ref{ch:open-obstructions}.

\subsubsection*{Comparison with the denominator algebra}

The recognition theorem is the source side of the Vol III Drinfeld
double construction.  The target presentation uses the generalized
Kac--Moody convention of Borcherds \cite{BorcherdsGKM88} and Kac
\cite{KAC:1}; the automorphic correction algebra is the
Gritsenko--Nikulin algebra of \cite[Sections~3--4,
Proposition~3.1]{GN}, and the product-lift data are those of
\cite[Theorem~2.1, \S5.1.1]{GNII} together with
Borcherds' automorphic product \cite{Borcherds95,B}.  The identity
$\mathrm{CoHA}(\C^3)=W_{1+\infty}$ belongs to the local toric case;
the present recognition criterion is at the level of the protected primitive Hall
algebra, before Drinfeld doubling; after the compact source datum is
supplied, it identifies that algebra with $\mathfrak g_{\Delta_5}$, not
with the centred or
vacuum-evaluated double.

Entries $(\mathrm L 1)$--$(\mathrm L 6)$ name the constructed
objects.  The finite HN system in
\S\ref{subsec:eight-finite-constructions} supplies the charge windows,
Hall correspondences, orientation transports, Pfaffian lines, Koszul
source coalgebras, primitive-recognition data, and their local
obstruction classes.
